%%% FIELD compact-transfer-statement
For every deterministic compact interval \([d,D]\subset(0,\infty)\), let \(v_\varepsilon=-q(\varepsilon)\).  There are constants \(\varepsilon_0=\varepsilon_0(d,D)>0\) and \(M=M(d,D)<\infty\) such that, for every \(0<\varepsilon\leq\varepsilon_0\) and every \((a_1,a_2,a_3)\in[d,D]^3\),
\[
\left|\varepsilon^{-1}H(\varepsilon a_1,\varepsilon a_2,\varepsilon a_3)-\frac{a_1a_2}{a_3}\right|
\leq \frac{M}{v_\varepsilon^2}.
\]
Thus the transfer is uniform on each fixed positive compact coefficient interval, including points on either cancellation plane \(a_1=a_3\) or \(a_2=a_3\).
%%% FIELD compact-transfer-proof
Fix \(0<d\leq D<\infty\), put \(d_0=\min\{d,1\}\), \(D_0=\max\{D,1\}\), and set
\[
\kappa=\max\{|\log d|,|\log D|\},
\qquad R=\log(D/d).
\tag{1}
\]
We first record bounds with explicit uniform content.  For \(w>0\), a change of variables in the normal tail gives
\[
\frac{\Phi(-w)}{\phi(w)}
=\int_0^\infty e^{-wt-t^2/2}\,dt
=\frac1w\int_0^\infty e^{-s}e^{-s^2/(2w^2)}\,ds.
\tag{2}
\]
Since \(e^{-y}<1\) for \(y>0\), while \(e^{-y}\geq1-y\) for \(y\geq0\), equation (2) and
\(\int_0^\infty s^2e^{-s}\,ds=2\) imply, for \(w>1\),
\[
\frac1w(1-w^{-2})
\leq \frac{\Phi(-w)}{\phi(w)}<\frac1w.
\tag{3}
\]
Because \(m(-w)=\phi(w)/\Phi(-w)\), (3) yields, for \(w\geq2\),
\[
w<m(-w)\leq\frac{w}{1-w^{-2}}.
\tag{4}
\]
Direct differentiation of \(m=\phi/\Phi\) gives
\[
m'(v)=-m(v)\{v+m(v)\}.
\tag{5}
\]
Using (4) in (5), with \(v=-w\), therefore gives
\[
|m'(-w)|=m(-w)\{m(-w)-w\}
\leq(1-w^{-2})^{-2}\leq\frac{16}{9}<2,
\qquad w\geq2.
\tag{6}
\]

Let \(v=v_\varepsilon=-q(\varepsilon)\).  Since \(\varepsilon\downarrow0\) implies \(v\to\infty\), we may choose \(\varepsilon_0\) so that, whenever \(0<\varepsilon\leq\varepsilon_0\),
\[
v\geq4,
\qquad e^{3v^2/8}>D_0,
\qquad v^2\geq12\kappa,
\qquad \varepsilon D_0<1.
\tag{7}
\]
For \(s>0\), differentiation and (4) show
\[
\frac{d}{ds}\log\Phi(-s)=-m(-s)<-s.
\]
Consequently,
\[
\log\frac{\Phi(-v/2)}{\Phi(-v)}
=\int_{v/2}^{v}m(-s)\,ds
>\int_{v/2}^{v}s\,ds
=\frac{3v^2}{8}.
\tag{8}
\]
As \(\Phi(-v)=\varepsilon\), (7)--(8) imply
\(\varepsilon a<\Phi(-v/2)\) for every \(a\in[d_0,D_0]\), and hence
\[
q(\varepsilon a)<-v/2.
\tag{9}
\]
The inverse-function rule and the definition of \(m\) give
\[
\frac{d}{dt}q(\varepsilon e^t)
=\frac{\varepsilon e^t}{\phi(q(\varepsilon e^t))}
=\frac1{m(q(\varepsilon e^t))}.
\tag{10}
\]
Integrating (10) from \(0\) to \(\log a\), and using (4) and (9) throughout that segment, proves the uniform quantile-shift bound
\[
q(\varepsilon a)=-v+\delta_\varepsilon(a),
\qquad
|\delta_\varepsilon(a)|\leq\frac{2\kappa}{v},
\qquad a\in[d,D].
\tag{11}
\]

Now take \((a_1,a_2,a_3)\in[d,D]^3\), and in \(\texttt{def:mixed-coefficient}\) put
\[
b=\varepsilon a_3,
\qquad x=\log(a_1/a_3),
\qquad z=\log(a_2/a_3).
\]
For \(s,t\in[0,1]\), the coefficients \(a_3e^{sx}\) and \(a_3e^{tz}\) are geometric interpolants of points in \([d,D]\), so they remain in \([d,D]\).  By (11), each of the three quantiles appearing in the mixed-coefficient integrand is \(-v\) plus an error of absolute value at most \(2\kappa/v\).  Its numerator is evaluated at
\[
h_{s,t}=q(be^{sx})+q(be^{tz})-q(b)
=-v+e_{s,t},
\qquad |e_{s,t}|\leq\frac{6\kappa}{v}.
\tag{12}
\]
Conditions (7) and (12) give \(-h_{s,t}\geq v/2\geq2\).  Equations (4), (6), and (9) therefore bound the absolute value of the integrand defining \(B\) by
\[
\frac{|m'(h_{s,t})|}
{m(q(be^{sx}))m(q(be^{tz}))}
\leq\frac{2}{(v/2)^2}=\frac8{v^2}.
\tag{13}
\]
After integration over the unit square,
\[
|B(\varepsilon a_1,\varepsilon a_2,\varepsilon a_3)|
\leq\frac8{v^2}.
\tag{14}
\]

The global exact factorization \(\texttt{thm:global-exact-factorization}\) now gives, with
\(A=\log(a_1/a_3)\log(a_2/a_3)\),
\[
\varepsilon^{-1}H(\varepsilon a_1,\varepsilon a_2,\varepsilon a_3)
=\frac{a_1a_2}{a_3}\exp\{A B(\varepsilon a_1,\varepsilon a_2,\varepsilon a_3)\}.
\tag{15}
\]
Here \(|A|\leq R^2\) and \(a_1a_2/a_3\leq D^2/d\).  Decreasing \(\varepsilon_0\), if necessary, so that \(8R^2/v^2\leq1\), and using \(|e^y-1|\leq e^{|y|}|y|\), equations (14)--(15) yield
\[
\left|\varepsilon^{-1}H(\varepsilon a_1,\varepsilon a_2,\varepsilon a_3)-\frac{a_1a_2}{a_3}\right|
\leq \frac{8eD^2R^2}{d\,v^2}.
\tag{16}
\]
Thus the assertion holds, for example, with \(M=1+8eD^2R^2/d\).  No step divided by \(A\), so (16) also holds on both exact cancellation planes.
%%% FIELD product-poisson-statement
If \(\tau_n=n\rho_n\to\tau\in[0,\infty)\), then the sixty-vector \(K\Rightarrow Z\), and jointly \(n\widehat\Delta^P_{\ell}\Rightarrow T_{\ell}\), \(n\widehat\Delta^E_{\ell}\Rightarrow T_{\ell}\), and \(n\Delta_{\ell}\to\tau(Q_{\ell}-c_{\ell4})\). Hence \(nL^{-1}\sum_{\ell=1}^{L}\widehat\Delta^E_{\ell}\Rightarrow L^{-1}\sum_{\ell=1}^{L}T_{\ell}\); at \(\tau=0\), this limit equals \(L^{-1}\sum_{\ell=1}^{L}\{\pi_{\ell3}/(2\pi_{\ell1}\pi_{\ell2})-1/(2\pi_{\ell4})\}\), which need not vanish under unequal shares.
%%% FIELD product-poisson-proof
Write \(N_{\ell j,n}=n_{\ell j}\), \(p_{\ell j,n}=r_{\ell j,n}\), and
\[
\lambda_{\ell j}=\pi_{\ell j}\tau c_{\ell j}.
\]
By \(\texttt{ass:bernoulli-role-sampling}\), \(K_{\ell j,n}\) has the binomial law with parameters \(N_{\ell j,n}\) and \(p_{\ell j,n}\).  By \(\texttt{ass:common-rarity}\), \(\texttt{ass:role-share-limits}\), \(\texttt{ass:coefficient-limits}\), and the displayed premise \(n\rho_n\to\tau\),
\[
N_{\ell j,n}p_{\ell j,n}
=\frac{N_{\ell j,n}}n(n\rho_n)c_{\ell j,n}
\longrightarrow\lambda_{\ell j}.
\tag{17}
\]
The rarity limit and coefficient limits imply \(p_{\ell j,n}=\rho_nc_{\ell j,n}\to0\), and therefore
\[
N_{\ell j,n}p_{\ell j,n}^2
=(N_{\ell j,n}p_{\ell j,n})p_{\ell j,n}\longrightarrow0.
\tag{18}
\]
Moreover, \(\texttt{ass:role-share-bounds}\) and the share limits imply
\(\underline\pi\leq\pi_{\ell j}\leq\overline\pi\); in particular, every \(\pi_{\ell j}\) is positive and \(N_{\ell j,n}\to\infty\).

For every fixed nonnegative integer \(k\), the exact binomial mass is
\[
\Pr(K_{\ell j,n}=k)
=\frac1{k!}\prod_{h=0}^{k-1}(N_{\ell j,n}-h)p_{\ell j,n}
\,(1-p_{\ell j,n})^{N_{\ell j,n}}(1-p_{\ell j,n})^{-k}.
\tag{19}
\]
Equations (17)--(18) give
\[
N_{\ell j,n}\log(1-p_{\ell j,n})
=-N_{\ell j,n}p_{\ell j,n}+O(N_{\ell j,n}p_{\ell j,n}^2)
\longrightarrow-\lambda_{\ell j}.
\tag{20}
\]
For \(k=0\), the product in (19) is empty; for \(k\geq1\), each of its \(k\) factors converges to \(\lambda_{\ell j}\).  Thus (19)--(20) prove
\[
\Pr(K_{\ell j,n}=k)
\longrightarrow e^{-\lambda_{\ell j}}\frac{\lambda_{\ell j}^k}{k!},
\tag{21}
\]
with the same conclusion when \(\lambda_{\ell j}=0\).

The disjoint-role independence assumption makes the sixty counts mutually independent.  Hence their joint mass at each point of \(\mathbb N^{4L}\) converges, by (21), to the product of the Poisson masses in \(\texttt{def:bounded-count-limit}\).  In addition, (17) and the fixed dimension give a finite constant \(C_0\) such that, for all sufficiently large \(n\),
\[
\mathbb E\!\left[\sum_{\ell=1}^{L}\sum_{j=1}^{4}K_{\ell j,n}\right]
=\sum_{\ell=1}^{L}\sum_{j=1}^{4}N_{\ell j,n}p_{\ell j,n}
\leq C_0.
\tag{22}
\]
Therefore
\[
\Pr\!\left(\max_{\ell,j}K_{\ell j,n}>M\right)
\leq\frac{C_0}{M+1}.
\tag{23}
\]
Pointwise convergence of the joint masses on the finite cube \(\{0,\ldots,M\}^{4L}\), followed by (23) and then \(M\to\infty\), proves directly that
\[
K\Rightarrow Z,
\tag{24}
\]
where the coordinates of \(Z\) are independent and have the stated Poisson means.  Indeed, the same Markov bound applies to \(Z\), because \(\mathbb E\sum_{\ell,j}Z_{\ell j}=\sum_{\ell,j}\lambda_{\ell j}<\infty\); hence the omitted mass outside the finite cube tends to zero for both sides.

Define the positive rescaled add-half cells
\[
a_{\ell j,n}=n\widetilde r_{\ell j}
=\frac{K_{\ell j,n}+1/2}{(N_{\ell j,n}+1)/n}.
\tag{25}
\]
The deterministic denominators in (25) converge to \(\pi_{\ell j}>0\).  On every fixed finite count cube, the map from the count vector to the vector in (25) therefore converges uniformly to
\[
k\longmapsto\left(\frac{k_{\ell j}+1/2}{\pi_{\ell j}}\right)_{\ell,j}.
\]
Combining this finite-cube uniform convergence with (23)--(24) proves, without any unbounded mapping step,
\[
(a_{\ell j,n})_{\ell,j}\Rightarrow
(W_{\ell j})_{\ell,j},
\qquad W_{\ell j}=\frac{Z_{\ell j}+1/2}{\pi_{\ell j}}.
\tag{26}
\]
Every coordinate in (25)--(26) is strictly positive, including when its count is zero.  The same finite-cube argument applied to the rational map gives jointly over \(1\leq\ell\leq L\),
\[
n\widehat\Delta^P_{\ell}
=\frac{a_{\ell1,n}a_{\ell2,n}}{a_{\ell3,n}}-a_{\ell4,n}
\Rightarrow
\frac{W_{\ell1}W_{\ell2}}{W_{\ell3}}-W_{\ell4}
=T_{\ell}.
\tag{27}
\]
The divisions in (27) are valid because the add-half numerator is at least \(1/2\).

We next compare the exact and multiplicative links.  Fix an integer \(M\geq0\), and define
\[
\mathcal F_{n,M}=\left\{\max_{\ell,j}K_{\ell j,n}\leq M\right\},
\qquad
d_M=\frac1{2(\overline\pi+1)},
\qquad
D_M=\frac{M+1/2}{\underline\pi}.
\tag{28}
\]
On \(\mathcal F_{n,M}\), the share bounds and (25) give, for every cell and every \(n\),
\[
0<d_M\leq a_{\ell j,n}\leq D_M<\infty.
\tag{29}
\]
Indeed, \(N_{\ell j,n}+1\leq n(\overline\pi+1)\) gives the lower bound, and \(N_{\ell j,n}+1\geq n\underline\pi\) gives the upper bound.  Apply \(\texttt{lem:fixed-compact-gaussian-tail-transfer}\) to this deterministic interval \([d_M,D_M]\), with \(\varepsilon=1/n\).  If \(v_n=-q(1/n)\), then \(v_n\to\infty\), and there is a finite constant \(C_M\), independent of \(n\), the pair, and the realized count vector in \(\mathcal F_{n,M}\), such that
\[
\max_{1\leq\ell\leq L}
\left|
nH(a_{\ell1,n}/n,a_{\ell2,n}/n,a_{\ell3,n}/n)
-\frac{a_{\ell1,n}a_{\ell2,n}}{a_{\ell3,n}}
\right|
\leq\frac{C_M}{v_n^2}
\quad\hbox{on }\mathcal F_{n,M}.
\tag{30}
\]
The arguments of \(H\) are exactly the add-half probabilities and hence lie in \((0,1)\), even for all-zero or all-full count cells.  By (23), for every \(\delta>0\), (30) implies
\[
\limsup_{n\to\infty}
\Pr\!\left(
\max_{1\leq\ell\leq L}
|n\widehat\Delta^E_{\ell}-n\widehat\Delta^P_{\ell}|>\delta
\right)
\leq\frac{C_0}{M+1}.
\tag{31}
\]
Letting \(M\to\infty\) proves
\[
\max_{1\leq\ell\leq L}
|n\widehat\Delta^E_{\ell}-n\widehat\Delta^P_{\ell}|
\longrightarrow_P0.
\tag{32}
\]
Equations (27) and (32), again checked first on finite count cubes and then using (23), prove the asserted joint exact-link statistic limit.

It remains to relate the observed-law functional to the causal target.  The finite collection of positive limits in \(\texttt{ass:coefficient-limits}\) permits fixed constants \(0<d_\star<D_\star<\infty\) such that
\[
c_{\ell j,n}\in[d_\star,D_\star]
\quad\hbox{for every }\ell,j\hbox{ and all sufficiently large }n.
\tag{33}
\]
Apply the fixed-compact transfer lemma to \([d_\star,D_\star]\), now with \(\varepsilon=\rho_n\).  Put \(u_n=-q(\rho_n)\), and define the proof intermediate
\[
Q_{\ell,n}=\frac{c_{\ell1,n}c_{\ell2,n}}{c_{\ell3,n}}.
\]
The rarity limit gives \(u_n\to\infty\), and the lemma yields a finite \(C_\star\) such that
\[
\max_{1\leq\ell\leq L}
\left|
\rho_n^{-1}H(\rho_nc_{\ell1,n},\rho_nc_{\ell2,n},\rho_nc_{\ell3,n})
-Q_{\ell,n}
\right|
\leq\frac{C_\star}{u_n^2}.
\tag{34}
\]
By \(\texttt{prop:published-identification-baseline}\), the causal contrast is identified by the observed-law map,
\[
\Delta_{\ell,n}
=H(r_{\ell1,n},r_{\ell2,n},r_{\ell3,n})-r_{\ell4,n}.
\tag{35}
\]
Substitution of common rarity into (35), followed by (34), gives
\[
n\Delta_{\ell,n}
=n\rho_n\{Q_{\ell,n}-c_{\ell4,n}\}
+O\!\left(\frac{n\rho_n}{u_n^2}\right).
\tag{36}
\]
The sequence \(n\rho_n\) is bounded, \(u_n\to\infty\), and coefficient convergence with \(c_{\ell3}>0\) gives \(Q_{\ell,n}\to Q_{\ell}\).  Hence (36) proves
\[
n\Delta_{\ell,n}\longrightarrow
\tau(Q_{\ell}-c_{\ell4}).
\tag{37}
\]
This establishes the claimed relationship between the causal target and the observed-law transfer; it is not imposed by definition.

Because \(L=15\) is fixed, summing the joint exact-link limits and dividing by \(L\) gives
\[
nL^{-1}\sum_{\ell=1}^{L}\widehat\Delta^E_{\ell}
\Rightarrow L^{-1}\sum_{\ell=1}^{L}T_{\ell}.
\tag{38}
\]
Finally, if \(\tau=0\), then every Poisson mean is zero and every \(Z_{\ell j}=0\) almost surely.  Therefore
\[
W_{\ell j}=\frac1{2\pi_{\ell j}},
\qquad
T_{\ell}
=\frac{\pi_{\ell3}}{2\pi_{\ell1}\pi_{\ell2}}
-\frac1{2\pi_{\ell4}},
\tag{39}
\]
which proves the displayed aggregate limit.  Positivity of all shares validates every denominator.  Unequal shares do not force (39) to be zero: whenever the share bounds permit two distinct positive values \(a\) and \(b\), shares satisfying \(\pi_{\ell1}=\pi_{\ell2}=\pi_{\ell3}=a\) and \(\pi_{\ell4}=b\neq a\) give \(T_{\ell}=1/(2a)-1/(2b)\neq0\).  Thus the all-zero limit is retained by the add-half rule and may exhibit the stated unequal-share displacement.
